% =============================================================================
% Chapter 16: Conclusion — ML Systems Lecture Slides
% =============================================================================
% This is the capstone lecture. It synthesizes all of Volume I: the twelve
% quantitative invariants, Conservation of Complexity, lighthouse model
% journeys, the system-is-the-model thesis, compound AI systems, and the
% bridge to Volume II.
%
% IMAGES NEEDED (SVGs converted to PDF):
%   - twelve-invariants-map.pdf    - iron-law-retrospective.pdf
%   - course-journey-timeline.pdf  - node-to-fleet.pdf
%   - conservation-of-complexity.pdf
% =============================================================================
\documentclass[aspectratio=169, 12pt]{beamer}
\usepackage{../../assets/beamerthememlsys}

\mlsyssetup{
  volume   = {Volume I},
  chapter  = {Chapter 16},
  logo = {../../assets/img/logo-mlsysbook.png},
  instlogo = {../../assets/img/logo-harvard.png},
  chaptertitle = {Conclusion},
}

% --- Fonts ---
\usepackage{fontspec}
\setsansfont{Helvetica Neue}[
  BoldFont={Helvetica Neue Bold},
  ItalicFont={Helvetica Neue Italic},
  BoldItalicFont={Helvetica Neue Bold Italic},
]
\setmonofont{JetBrains Mono}[Scale=0.85]

% --- Packages ---
\usepackage{booktabs}
\usepackage{amsmath}

% --- Image paths ---
\graphicspath{
  {images/}
}

% --- Chapter-specific macros ---
\newcommand{\DAM}{D\raisebox{0.08em}{\tiny$\bullet$}A\raisebox{0.08em}{\tiny$\bullet$}M}

% --- Helper: safe image include ---
\newcommand{\safeimg}[2][width=\textwidth,keepaspectratio]{%
  \IfFileExists{images/#2}{\includegraphics[#1]{#2}}{%
    \IfFileExists{#2}{\includegraphics[#1]{#2}}{%
      \fbox{\parbox[c][2.5cm][c]{0.85\linewidth}{\centering\footnotesize\textcolor{midgray}{[Missing image]}}}%
    }%
  }%
}

% --- Section count for navigation (must match actual \section{} count) ---
\setcounter{mlsystotalsections}{6}

\title{Conclusion}
\author{Vijay Janapa Reddi}
\institute{Harvard University}
\date{}

\begin{document}

% =============================================================================
% TITLE SLIDE
% =============================================================================
\mlsystitle{Conclusion}{The System Is the Model}{cover_conclusion.png}

% =============================================================================
% LEARNING OBJECTIVES
% =============================================================================
\begin{frame}{Learning Objectives}
\note{[2 min] This is the synthesis lecture. Students already know everything
from 15 chapters. Today we connect the dots. Emphasize that no new content
is introduced --- only integration and retrieval.
Ask: ``What is the single most important thing you learned this semester?''}

\small
\begin{enumerate}
  \item Synthesize the \textbf{twelve quantitative invariants} into an integrated framework
  \item Apply the \textbf{Conservation of Complexity} meta-principle to trace trade-offs
  \item Trace how \textbf{Lighthouse Models} reveal constraint propagation across the stack
  \item Evaluate the \textbf{Iron Law} as a prescriptive optimizer, not just a diagnostic
  \item Connect technical foundations, scale engineering, and production reality
  \item Identify how single-node principles extend to \textbf{fleet-scale systems}
\end{enumerate}

\end{frame}

% =============================================================================
\section{The System Is the Model}
% =============================================================================

\begin{frame}{The System Is the Model}
\note{[3 min] Open with the mobile deployment failure from the chapter.
The key insight: no single team caused the failure. The interaction
between quantization, firmware, and preprocessing did.
Ask: ``Has anyone seen a bug that no single team owned?''}

\footnotesize
\begin{columns}[T]
  \begin{column}{0.55\textwidth}
    A model that produces perfect predictions is useless if:
    \begin{itemize}\setlength\itemsep{0pt}
      \item It receives \textbf{corrupted inputs}
      \item It cannot be \textbf{deployed reliably}
      \item It \textbf{drifts} without anyone noticing
    \end{itemize}

    \vspace{0.15cm}
    \begin{mlsyscard}{crimson}
    {\scriptsize The \emph{true model} = data pipeline + training infra + serving system + monitoring loop.}
    \end{mlsyscard}

    \vspace{0.1cm}
    \mlsysconcept{Thesis:}{Optimize the system, the model improves.}
  \end{column}
  \begin{column}{0.42\textwidth}
    \renewcommand{\arraystretch}{1.1}
    {\scriptsize
    \begin{tabular}{@{}ll@{}}
      \toprule
      \textbf{Component} & \textbf{Failure Mode} \\
      \midrule
      Data pipeline & Schema drift \\
      Training & Stale features \\
      Compression & Numerical skew \\
      Serving & Latency violation \\
      Monitoring & Silent degradation \\
      \bottomrule
    \end{tabular}
    }
  \end{column}
\end{columns}

\end{frame}

% --- ACTIVE LEARNING 1: Predict ---
\begin{frame}{Predict: How Many Invariants?}
\note{[2 min] Prediction exercise before revealing the twelve invariants.
Give students 60 seconds to write. Most will guess 3--5.
The reveal (12) is surprising and motivates the next section.}

\centering
\vspace{1.0cm}
{\Large\bfseries Think--Write--Share}

\vspace{0.6cm}
{\large How many \textbf{quantitative invariants} (physical laws,\\
not best practices) have we encountered across 15 chapters?}

\vspace{0.6cm}
{\normalsize Write a number and name as many as you can. \textcolor{midgray}{(60 seconds)}}

\vspace{0.6cm}
{\small\textcolor{lightgray}{Turn to a neighbor and compare your lists.}}

\end{frame}

% =============================================================================
\section{Twelve Invariants}
% =============================================================================

\begin{frame}{Twelve Invariants, Four Phases}
\note{[3 min] Reveal the full map. Walk through each quadrant briefly.
The key insight: these are from physics, information theory, and statistics
--- not trends or best practices. They will still be valid in 10 years.
Ask: ``Which quadrant surprised you the most?''}

% --- Layout: FULL-WIDTH diagram ---
\centering
\safeimg[width=0.92\textwidth,height=4.8cm,keepaspectratio]{twelve-invariants-map.pdf}

\vspace{0.1cm}
\mlsysinsight{Key:}{Four phases form a cycle. Deploy feeds back to Foundations through drift and retraining.}

\end{frame}

\begin{frame}{Foundations \& Build (Invariants 1--4)}
\note{[2 min] Cover the first four invariants quickly. Students know these
well from the first half of the course. Emphasize that Data as Code and
Data Gravity are ``where complexity originates'' while Iron Law and Silicon
Contract are ``how complexity becomes computation.''}

\footnotesize
\begin{columns}[T]
  \begin{column}{0.48\textwidth}
    \textcolor{computestroke}{\textbf{I: Foundations}} --- Where complexity originates

    \vspace{0.15cm}
    \textbf{1. Data as Code}\\
    {\scriptsize System behavior $\approx f(\text{Data})$. Changing data changes the program.}

    \vspace{0.1cm}
    \textbf{2. Data Gravity}\\
    {\scriptsize $C_{move}(\text{Data}) \gg C_{move}(\text{Compute})$. Move compute to data.}

    \vspace{0.15cm}
    \mlsysconcept{Lighthouse:}{DLRM is Data Gravity embodied --- TB-scale embeddings anchor the architecture.}
  \end{column}
  \begin{column}{0.48\textwidth}
    \textcolor{datastroke}{\textbf{II: Build}} --- How complexity becomes computation

    \vspace{0.15cm}
    \textbf{3. Iron Law}\\
    {\scriptsize $T = D_{vol}/BW + O/(R_p \cdot \eta) + L_{lat}$. Three levers.}

    \vspace{0.1cm}
    \textbf{4. Silicon Contract}\\
    {\scriptsize Every architecture bets on which hardware resource it saturates.}

    \vspace{0.15cm}
    \mlsysconcept{Lighthouse:}{ResNet-50 is compute-bound; GPT-2/Llama is bandwidth-bound. Same equation, different dominant term.}
  \end{column}
\end{columns}

\end{frame}

\begin{frame}{Optimize \& Deploy (Invariants 5--12)}
\note{[3 min] Cover the remaining eight invariants. Emphasize the tightly
coupled diagnostic chain in Optimize and the ``reality defeats assumptions''
theme in Deploy. The deployment invariants detect failures the first eight
cannot prevent.}

\footnotesize
\begin{columns}[T]
  \begin{column}{0.48\textwidth}
    \textcolor{routingstroke}{\textbf{III: Optimize}} --- Constraints shape trade-offs

    \vspace{0.1cm}
    \textbf{5. Pareto Frontier} --- No free improvements\\
    \textbf{6. Arithmetic Intensity} --- Diagnoses bottleneck\\
    \textbf{7. Energy-Movement} --- Data locality dominates\\
    \textbf{8. Amdahl's Law} --- Serial fraction caps gains

    \vspace{0.15cm}
    {\scriptsize MobileNetV2 navigates all four: depthwise separable convolutions reshape Pareto; INT8 exploits arithmetic intensity; energy savings respect Energy-Movement; Amdahl's explains preprocessing bottleneck.}
  \end{column}
  \begin{column}{0.48\textwidth}
    \textcolor{errorstroke}{\textbf{IV: Deploy}} --- Reality defeats assumptions

    \vspace{0.1cm}
    \textbf{9. Verification Gap} --- Statistical testing only\\
    \textbf{10. Statistical Drift} --- Accuracy erodes over time\\
    \textbf{11. Training-Serving Skew} --- Subtle divergences\\
    \textbf{12. Latency Budget} --- P99 is the hard constraint

    \vspace{0.15cm}
    {\scriptsize DLRM loses revenue if Training-Serving Skew corrupts features. GPT-2/Llama must respect Latency Budget via continuous batching. A chatbot that responds in 10 s is a chatbot nobody uses.}
  \end{column}
\end{columns}

\end{frame}

% =============================================================================
\section{Conservation of Complexity}
% =============================================================================

\mlsysfocus{Conservation of Complexity}{%
Complexity cannot be destroyed --- only moved\\[0.2cm]
between \textcolor{yellow}{Data}, \textcolor{cyan}{Algorithm}, and \textcolor{orange}{Machine}.%
}

\begin{frame}{Conservation of Complexity in Action}
\note{[3 min] The meta-principle that unifies all twelve invariants.
Walk through both examples. Quantization: you reduced BW demand but
increased monitoring burden. Abstraction: you simplified the interface
but hid resource consumption. The Conservation of Complexity is analogous
to conservation of energy in physics.
Ask: ``Can you name an optimization where the hidden cost surprised you?''}

% --- Layout: FULL-WIDTH diagram ---
\centering
\safeimg[width=0.92\textwidth,height=4.8cm,keepaspectratio]{conservation-of-complexity.pdf}

\vspace{0.1cm}
\mlsysalert{Rule:}{Every optimization is a complexity \emph{transfer}, not a complexity \emph{elimination}.}

\end{frame}

\begin{frame}{One Decision, Six Invariants}
\note{[3 min] Trace the quantization decision through the invariant web.
This is the key ``aha'' moment: a single engineering decision (FP16 to INT8)
ripples through Pareto, Silicon Contract, Arithmetic Intensity,
Energy-Movement, Latency Budget, and Training-Serving Skew.
Walk through the diagram's radiating spokes.}

% --- Layout: FULL-WIDTH diagram ---
\centering
\safeimg[width=0.92\textwidth,height=4.5cm,keepaspectratio]{iron-law-retrospective.pdf}

\vspace{0.1cm}
{\small\textbf{Profile first. Optimize second.} Identify the dominant Iron Law term before investing effort.}

\end{frame}

% --- ACTIVE LEARNING 2: Exercise ---
\begin{frame}{Your Turn: Trace the Invariants}
\note{[3 min] Give students 90 seconds. This is a synthesis exercise ---
they must connect multiple invariants to a single decision.
Answer: Pruning pulls the Compute term (fewer FLOPs) AND the Data Movement
term (fewer parameters to load). It navigates the Pareto Frontier
(accuracy vs speed). The Silicon Contract determines whether the hardware
can exploit sparsity. Training-Serving Skew must be checked because
pruned weights may behave differently at inference.}

\small
\begin{columns}[T]
  \begin{column}{0.55\textwidth}
    {\normalsize\bfseries Trace the Invariants}

    \vspace{0.2cm}
    A colleague proposes \textbf{pruning} a ResNet-50 model to 50\% sparsity for mobile deployment.

    \vspace{0.15cm}
    \textbf{Which invariants are affected?}
    \begin{itemize}\setlength\itemsep{0pt}
      \item Which Iron Law terms change?
      \item What does the Pareto Frontier predict?
      \item What must you check for Training-Serving Skew?
    \end{itemize}

    {\footnotesize\textcolor{midgray}{(90 seconds --- then discuss with a neighbor)}}
  \end{column}
  \begin{column}{0.42\textwidth}
    \pause
    \begin{mlsyscard}{datastroke}
    \textbf{Solution:}\\[0.1cm]
    {\footnotesize
    \textbf{Iron Law:} Compute term (fewer FLOPs) + Data term (fewer params to load)\\[0.05cm]
    \textbf{Pareto:} Trading capacity for speed; accuracy drops\\[0.05cm]
    \textbf{Silicon Contract:} Does hardware exploit sparsity?\\[0.05cm]
    \textbf{Skew:} Pruned inference path must match training behavior
    }
    \end{mlsyscard}
  \end{column}
\end{columns}

\end{frame}

% =============================================================================
\section{Lighthouse Journeys}
% =============================================================================

\begin{frame}{Five Lighthouses, Five Bottlenecks}
\note{[2 min] Quick retrieval: students should be able to name all five
Lighthouses and their bottlenecks by now. This is a callback to Ch 1.
Walk through the timeline diagram showing where each lighthouse appeared.}

% --- Layout: FULL-WIDTH diagram ---
\centering
\safeimg[width=0.92\textwidth,height=3.5cm,keepaspectratio]{course-journey-timeline.pdf}

\vspace{0.05cm}
\scriptsize
\renewcommand{\arraystretch}{0.95}
\begin{tabular}{@{}llll@{}}
  \toprule
  \textbf{Lighthouse} & \textbf{Bottleneck} & \textbf{Iron Law} & \textbf{Key Insight} \\
  \midrule
  ResNet-50    & \textcolor{computestroke}{Compute}  & $O/R_p$ & Batch size transforms bound \\
  GPT-2/Llama  & \textcolor{datastroke}{Memory BW}  & $D/BW$  & KV-cache dominates serving \\
  MobileNetV2  & \textcolor{routingstroke}{Efficiency} & All    & Pareto across 7 phases \\
  DLRM         & \textcolor{errorstroke}{Capacity}   & $D/BW$  & TB embeddings anchor arch. \\
  KWS          & Every resource & All     & Sub-MB, microwatt budget \\
  \bottomrule
\end{tabular}

\end{frame}

\begin{frame}{The Cost of a Token}
\note{[3 min] A worked example applying the Iron Law and Arithmetic
Intensity to LLM serving. The key reveal: memory time is 40x larger
than compute time. An engineer who optimizes compute kernels without
addressing memory wastes 100\% of their effort.
Ask: ``What are two ways to fix this?'' (Batching, quantization)}

\footnotesize
\textbf{Llama-2-70B inference on NVIDIA H100}

\vspace{0.05cm}
\renewcommand{\arraystretch}{1.05}
{\scriptsize
\begin{tabular}{@{}lllr@{}}
  \toprule
  \textbf{Quantity} & \textbf{Value} & \textbf{Calculation} & \textbf{Time} \\
  \midrule
  Model size ($D_{vol}$) & 70B $\times$ 2 B = 140 GB & & \\
  Compute ($O$) & 2 $\times$ 70B = 140 GFLOPs & & \\
  \midrule
  \textcolor{datastroke}{Memory time} & 140 GB / 3.35 TB/s & $D_{vol}/BW$ & \textbf{41.8 ms} \\
  \textcolor{computestroke}{Compute time} & 140 GF / 989 TF & $O/R_p$ & \textbf{0.14 ms} \\
  \bottomrule
\end{tabular}
}

\vspace{0.1cm}
\begin{mlsyscard}{errorstroke}
{\scriptsize\textbf{Ratio:} Memory time is \textbf{$\approx$300$\times$} larger than compute time. Optimizing compute without addressing memory wastes 100\% of effort.}
\end{mlsyscard}

\vspace{0.05cm}
{\scriptsize\textbf{Fix:} Batch users to reuse $D_{vol}$ (increase arith.\ intensity) \textbf{or} quantize to INT4 (reduce $D_{vol}$).}

\end{frame}

% --- ACTIVE LEARNING 3: Discussion ---
\begin{frame}{Discussion: What Would You Change?}
\note{[3 min] Turn-and-talk. This is a retrospective discussion.
Students reflect on what they would do differently knowing what they know
now. Cold-call 2--3 pairs. The point: systems thinking changes how you
approach every single-component decision.}

\centering
\vspace{0.8cm}
{\large\bfseries Turn and Talk \textcolor{midgray}{(90 seconds)}}

\vspace{0.6cm}
{\large Knowing what you know now about the twelve invariants,\\[0.2cm]
\alert{what would you do differently} in a past ML project?}

\vspace{0.6cm}
{\normalsize Name the invariant you wish you had applied.}

\vspace{0.6cm}
{\small\textcolor{lightgray}{Be specific: which decision, which invariant, which outcome.}}

\end{frame}

% =============================================================================
\section{Future Directions}
% =============================================================================

\begin{frame}{Four Deployment Paradigms}
\note{[2 min] Quick overview of the four paradigms. Each foregrounds
different invariants but none escapes them. The deployment spectrum from
Ch 1 and Ch 2 recurs here at the end. Same physics, different dominant terms.}

\scriptsize
\renewcommand{\arraystretch}{1.0}
\begin{tabular}{@{}lp{2.8cm}p{2.8cm}l@{}}
  \toprule
  \textbf{Paradigm} & \textbf{Dominant Invariants} & \textbf{Key Challenge} & \textbf{Lighthouse} \\
  \midrule
  \textbf{Cloud}       & Iron Law, Silicon Contract & Throughput at scale         & ResNet, DLRM \\
  \textbf{Edge/Mobile}  & Energy-Movement, Pareto   & Power \& latency budgets   & MobileNetV2 \\
  \textbf{Gen.\ AI} & Arith.\ Intensity, Latency & Memory-bound decoding & GPT-2/Llama \\
  \textbf{TinyML}       & All twelve simultaneously   & Every byte, every mW       & KWS \\
  \bottomrule
\end{tabular}

\vspace{0.15cm}
\begin{mlsyscard}{crimson}
{\scriptsize What unites these paradigms is not their hardware but their \textbf{physics}: the same invariants govern all of them.}
\end{mlsyscard}

\vspace{0.05cm}
{\footnotesize\mlsysinsight{Pattern:}{Constraints at the edge catalyze innovation that benefits all tiers.}}

\end{frame}

\begin{frame}{Compound AI Systems}
\note{[3 min] Introduce the compound AI systems concept from the chapter.
Key insight: instead of one monolithic model, chain specialized components.
Each can be independently optimized, monitored, and debugged. The Pareto
Frontier predicts this trade-off; Conservation of Complexity demands it.
Ask: ``Why not just build a bigger model?''}

\footnotesize
\begin{columns}[T]
  \begin{column}{0.55\textwidth}
    \textbf{Beyond monolithic models:}

    \vspace{0.1cm}
    Chain specialized components rather than scaling one model:
    \begin{itemize}\setlength\itemsep{0pt}
      \item \textbf{Retriever} finds relevant information
      \item \textbf{Reasoner} processes it
      \item \textbf{Verifier} checks the output
    \end{itemize}

    \vspace{0.1cm}
    \begin{mlsyscard}{computestroke}
    {\scriptsize Each component has its own \textbf{Silicon Contract} and \textbf{Arith.\ Intensity} profile. Interfaces create natural \textbf{monitoring points}.}
    \end{mlsyscard}
  \end{column}
  \begin{column}{0.42\textwidth}
    \renewcommand{\arraystretch}{1.1}
    {\scriptsize
    \begin{tabular}{@{}lll@{}}
      \toprule
      & \textbf{Monolithic} & \textbf{Compound} \\
      \midrule
      Update    & Retrain all   & Per-component \\
      Debug     & Black box     & Inspect steps \\
      Monitor   & End-to-end    & Per-interface \\
      Optimize  & One profile   & Per-component \\
      \bottomrule
    \end{tabular}
    }
  \end{column}
\end{columns}
\mlsyscite{Zaharia et al., ``The Shift to Compound AI Systems,'' 2024}

\end{frame}

% =============================================================================
\section{Wrap-Up}
% =============================================================================

\begin{frame}{Fallacies}
\note{[2 min] Three fallacies from the chapter. Each assumes that
optimizing one dimension suffices.}

\small
\textbf{Fallacy:} \textit{Systems engineering complexity disappears with better tools.}\\
{\footnotesize Tools abstract complexity; they do not eliminate it. The Conservation of Complexity guarantees that simplifying one interface pushes complexity to another.}

\vspace{0.15cm}
\textbf{Fallacy:} \textit{Mastering individual components equals mastering the system.}\\
{\footnotesize A data schema change cascades through training, breaks quantization assumptions, and triggers silent accuracy degradation. Integration complexity exceeds the sum of component complexities.}

\vspace{0.15cm}
\textbf{Fallacy:} \textit{A single accuracy metric captures model quality.}\\
{\footnotesize The Pareto Frontier: accuracy, latency, throughput, memory, energy, fairness, cost. A 95\% acc model at 500 ms may be worse than 93\% at 50 ms when P99 is the hard constraint.}

\end{frame}

\begin{frame}{Pitfalls}
\note{[2 min] Two operational pitfalls. Both share a common root:
the temptation to reduce a system to its parts.}

\small
\textbf{Pitfall:} \textit{Optimizing a single invariant while ignoring the Conservation of Complexity.}\\
{\footnotesize When an optimization reduces latency by 50\%, ask where the cost went. Quantization shifts load to accuracy monitoring. Caching trades memory capacity for speed. Engineers who celebrate gains in one metric without tracing compensating costs build systems that fail in unexpected ways.}

\vspace{0.2cm}
\textbf{Pitfall:} \textit{Optimizing a single pipeline stage without profiling the full system.}\\
{\footnotesize Amdahl's Law: optimizing inference by 10$\times$ yields only 1.1$\times$ system speedup if data loading accounts for 90\% of latency. The Iron Law decomposes execution time precisely so engineers can identify the dominant term \emph{before} investing effort.}

\end{frame}

% --- RETRIEVAL PRACTICE ---
\begin{frame}{What Were the Key Ideas?}
\note{[2 min] Retrieval practice. Students write 90 seconds, no notes.
Walk around the room. This is the most important active learning moment
in the capstone lecture --- students must synthesize the entire course.}

\centering
\vspace{1.5cm}
{\Large\bfseries Close your notes.}

\vspace{0.8cm}
{\large Write down the \textbf{4 most important ideas} from this entire course.}

\vspace{0.8cm}
{\normalsize\textcolor{midgray}{90 seconds --- no peeking.}}

\end{frame}

\begin{frame}{Key Takeaways}
\note{[3 min] Reveal. Walk through each bullet. These are the seven
takeaways from the entire course, not just this lecture. Emphasize the
last bullet: the node is mastered, the fleet awaits.}

\scriptsize
\begin{itemize}\setlength\itemsep{0pt}
  \item \textbf{Twelve invariants define ML systems engineering}: From Data as Code through Latency Budget --- physics, not trends.
  \item \textbf{Conservation of Complexity unifies all twelve}: Complexity cannot be destroyed; only moved between Data, Algorithm, and Machine.
  \item \textbf{The system is the model}: Data pipeline + training infra + serving + monitoring = the true model.
  \item \textbf{The Iron Law is prescriptive}: Profile first, optimize the dominant term. One decision ripples through six invariants.
  \item \textbf{Every deployment context stresses different invariants}: Same physics, different dominant terms.
  \item \textbf{Technical excellence + ethical commitment}: Verification Gap and drift invariants apply equally to fairness.
  \item \textbf{Mastering the node prepares you for the fleet}: Principles scale to the Warehouse-Scale Computer.
\end{itemize}

\end{frame}

\begin{frame}{References}
\note{[1 min] Point students to the canonical references. Hennessy \&
Patterson's ``New Golden Age'' paper is the intellectual foundation.}

\small
\mlsysref{H\&P19}{Hennessy \& Patterson. ``A New Golden Age for Computer Architecture.'' CACM, 2019.}
\mlsysref{Sculley+15}{Sculley et al. ``Hidden Technical Debt in ML Systems.'' NeurIPS, 2015.}
\mlsysref{Karpathy17}{A. Karpathy. ``Software 2.0.'' 2017.}
\mlsysref{Zaharia+24}{Zaharia et al. ``The Shift to Compound AI Systems.'' 2024.}
\mlsysref{Barroso+19}{Barroso et al. \textit{The Datacenter as a Computer}. 3rd ed. 2019.}
\mlsysref{Sutton19}{R. Sutton. ``The Bitter Lesson.'' 2019.}

\end{frame}

\begin{frame}{Next: Volume II --- Orchestrating the Fleet}
\note{[2 min] The closing forward hook. This is the bridge to Volume II.
Emphasize: same physics, new variables. The unit is mastered; the
collective awaits. End with: ``Go build it well.''}

% --- Layout: FULL-WIDTH diagram ---
\centering
\safeimg[width=0.88\textwidth,height=3.5cm,keepaspectratio]{node-to-fleet.pdf}

\vspace{0.05cm}
{\small\textbf{Same physics. New variables.} The Iron Law spans racks and zones.}

\vspace{0.1cm}
\begin{mlsyscard}{crimson}
\centering
{\small\textit{The future of intelligence is not a destiny we will merely witness.
It is a system we must engineer.} \textbf{Go build it well.}}
\end{mlsyscard}

\end{frame}

\end{document}
